% =====================================================================
\section{What the model predicts that is measurable}\label{sec:pred}
All of the work is done \emph{in silico}, and the first question put to it is what can be measured from here. Below is everything the model asserts about quantities accessible to experiment: a number, a direction, and what exactly refutes the assertion. The predictions are ordered by the strength of the test, not by the importance of the mechanism.

\begin{note}
A caveat about units, and it is obligatory. What is comparable is \emph{directions, orders of magnitude and ratios}, not absolute values. The parameters of the redox core are taken from the kinetics of cell lines and not of PV cells; the quantities of the sleep switch are expressed in millivolts of its own calibration. Where an absolute number stands below, it is quoted as the size of an effect inside the model, and it may not be carried over to experiment without recalculation.
\end{note}

\begin{enumerate}[leftmargin=1.5em,itemsep=4pt]
\item \textbf{Non-monotonicity of damage in vulnerability: the rescued band.} The population lives in three phases, not two. Between the healthy core and the collapsed one lies a band of cells which are \emph{more} vulnerable than their collapsed neighbours and are nonetheless intact: a fraction $0.00732\pm0.00042$ (seventeen standard errors from zero), and about $8.5\,\%$ of the population on a fine grid. With the slow defence frozen the band becomes exactly zero; under priming it widens $1.82$-fold. \emph{To measure}: damage to PV cells must not be a monotone function of their initial vulnerability --- the tissue must contain a subgroup of cells with high vulnerability markers and preserved parvalbumin, adjacent to damaged ones. Proxies for vulnerability: integrity of the perineuronal net, mitochondrial markers, indices of the oxidative load on the cell. \emph{Refuted by}: a monotone ``more vulnerable, more damaged'' relation at single-cell level. This is the most distinctive prediction of the work: it is counter-intuitive, follows from no averaged picture, and requires single-cell rather than tissue resolution.

\item \textbf{A precursor in sleep, and its source.} Sleep deteriorates \emph{before} the collapse threshold, and what does it is the slow memory of the defence rather than slow noise in the load: when the slow memory is taken away the height of the precursor falls in 107 sets of 174, $p=0.0015$. The height in the last day before a collapse is $1.061$ of the set's own healthy level, against an analytic bracket for the switch of $1.0500$--$1.1280$. \emph{To measure}: a lengthening of waking of the order of a few per cent in the days preceding an episode. Data of this kind are plentiful --- actigraphy and polysomnography in clinical high-risk cohorts and in schizophrenia, where sleep disturbance before an exacerbation is described. \emph{Refuted by}: the absence of a systematic shift before an episode; or by the presence of a shift that does not disappear under an intervention suppressing the slow link of the defence. The second is the stronger test but requires an animal model.

\item \textbf{The ceiling of passive repair, and the affine duration law.} The duration of an episode is set by the \emph{memory} of the defence and not by its strength: within a set, $99.8\,\%$ of the variance of the duration is explained by the course of the load and $0.2\,\%$ by repair, the decomposition being exact. The law is affine and does not saturate, duration $=1.28+0.97\,\tau_{\rm d}$ days, and a saturating form with a common offset is rejected. Hence: \emph{passive strengthening of repair does not attain the clinical target} --- missing by a factor of 2 to 10, since the required memory is about 12\,h whereas every measured value for the second stage lies in 24--88\,h. \emph{To measure}: interventions that lengthen or shorten the lifetime of the protective protein must shift the episode duration \emph{linearly}, without saturation; interventions that raise the strength of repair at unchanged memory must not shift it at all. \emph{Refuted by}: saturation of the response at large $\tau_{\rm d}$, or an appreciable shortening of episodes from strengthened repair at unchanged memory.

\item \textbf{Priming works through the top of the window, not through the speed of repair.} A preliminary mild load shortens the subsequent episode: the one-sided 95\,\% bound on the median duration is $1.51$ days, inside the clinical two. The mechanism is named: priming raises the \emph{top} of the window rather than accelerating repair. Quantitatively, accelerating the first stage of induction gives $-1$ to $-5\,\%$ of the duration, whereas strengthening it gives $-52\,\%$. \emph{To measure}: a hormesis protocol --- a mild oxidative exposure before the main one must shorten the episode. The two mechanisms are separated by an intervention that accelerates induction against one that raises its ceiling; the model predicts that only the second works. \emph{Refuted by}: equal efficacy of the two interventions, or an advantage for acceleration.

\item \textbf{The capacity loop: episodes rarer but longer.} Per-cell loss of defence capacity does not make collapses more frequent but makes them longer: duration $13.1\to25.6$ days with the rate falling $0.042\to0.033$ per set-day at risk, while episodicity is preserved at $0.75$ of the open model. Progression exceeds the capacity-matched static by $+0.0231$ at an MDE of $0.0096$, replicated on three independent noise realisations. \emph{To measure}: in longitudinal cohorts, as damage accumulates the episodes should become \emph{rarer and longer}, not more frequent. This is testable on data about the number and duration of exacerbations that are already being collected. \emph{Refuted by}: a growing rate of exacerbations under progression. Caveat: the absolute value of $25.6$ days lies far above the clinical two, so the working point of the loop is overstated; what is comparable here is the direction and the ratio.

\item \textbf{Hysteresis: retention depends on the depth to which the load is removed.} The broken state is retained after the load is removed, and the retention grows with the softness of the removal monotonically ($\rho=+1.000$, permutation $p=0.0083$): under the softest removal it is $0.561$ against $0.146$ for the open model. \emph{To measure}: the probability of relapse after treatment is withdrawn should depend monotonically on the completeness of the withdrawal rather than being constant. Testable in discontinuation studies, though ``completeness of removal'' needs an operationalisation of its own.

\item \textbf{A bursty load does not strike harder than a smooth one.} With the mean and the variance pinned, the collapse rate \emph{falls} by $40\,\%$ as the correlation time of the load grows. Rare strong pushes are no more dangerous than a steady rise --- a hypothesis we tested and refuted. \emph{To measure}: stress protocols with identical mean and variance but different temporal structure; clustered stress should give \emph{fewer} episodes than uniform stress. Testable in an animal model.

\item \textbf{A serotonin surge is a volume knob, not a trigger.} A surge is followed by $26\,\%$ of collapses, against an analytic expectation of $43\,\%$ and a pre-registered threshold of $80\,\%$; and no property of the episode changes. \emph{To measure}: serotonergic provocation should not reliably launch an episode and should not change its duration or depth.

\item \textbf{The asymmetry of the defence sets its mean level.} If the rise and decay constants of the protective protein differ by a factor $r$, the mean level of the defence stands at $w(r)=fr/(fr+1-f)$ with $f=0.705$, the measured fraction of waking. A scan of $r=0.33$--$3$ predicts $w=0.441$--$0.878$, and the measured levels of the defence follow that ordering exactly ($\rho=+1$, $p=0.0083$). The formula contains not one fitted number. \emph{To measure}: the two constants of the protective protein separately, in the cell. The model then predicts the mean level of the defence with no free parameters.

\item \textbf{Two negative predictions.} \emph{A loop through the common load cannot give chronicity and episodicity at once}: irreversibility engages at a loop strength of $1.11$ and the episodes die out at $0.18$ --- a sixfold structural gap, robust to a softer protocol. Consequence for experiment: if patients reliably show \emph{both} irreversible accumulation \emph{and} preserved episodicity, the mechanism does not lie in a loop through the common load. \emph{Breakdown of the rhythm as a locking staircase is unattainable at physiological coupling}: the first step requires $k\!\cdot\!D=4.77$\,mV while identically $k\!\cdot\!D\le k$, so at the working $k=2$\,mV it is not reached for any damage. The model thus predicts that sleep cycles of 48 and 72\,h should not occur in schizophrenia --- which agrees with their not being described. The prediction is weak, but it is the only place in the work where a disagreement between model and observation turned out to be in the model's favour.
\end{enumerate}

\textbf{What the model does \emph{not} predict.} Quantities of its internal construction --- the ratio of the window's boundaries, the positions of the folds in millivolts, the position of the load within the window, the number of cells in a set --- are not predictions: they are either chosen or auxiliary. Dopamine and the positive symptoms are outside the model by design. And, most importantly, the model describes the onset and the ending of an \emph{episode}, not the course of the disease, so any claim on the scale of years lies outside its domain of applicability.
